\chapter{पाचन और पोषक तत्व}\label{ch:g7:digestion-and-nutrients}

\Cref{ch:g4:journey-of-food} ने रोटी के कौर के पीछे-पीछे नली की सैर की
थी और वादा किया था कि ``\omterm{def:g4:journey-of-food:digestion}{पाचक रस} उसे टुकड़ा-टुकड़ा कर देते हैं''। इस साल
हम कारख़ाने के दरवाज़े खोलते हैं: वे रस ठीक-ठीक करते क्या हैं, वे
कौन-से \omterm{def:g7:digestion-and-nutrients:families}{पोषक तत्व} बनाते हैं, और \omterm{def:g4:journey-of-food:tract}{छोटी आँत} की दीवार शरीर की सबसे बड़ी
सरहद-पार आवाजाही कैसे सँभालती है? रोटी का कौर लौट आया है --- और इस बार
बेहतर औज़ारों के नीचे।

\section{रस, काम पर}

\begin{proposition}[पाचन रासायनिक तोड़-फोड़ है]\label{prop:g7:digestion-and-nutrients:chemical}
चबाना और मथना भोजन को बस \emph{छोटा} करते हैं; असली तोड़-फोड़ रासायनिक
है। \emph{पाचक रसों}\index{पाचक रस} में --- \omterm{def:g4:journey-of-food:tract}{लार}, \omterm{def:g4:journey-of-food:tract}{आमाशय} के रस, और \omterm{def:g4:journey-of-food:tract}{छोटी
आँत} में उँडेले जाने वाले रसों में --- ऐसे काम करने वाले पदार्थ होते हैं
जो भोजन के बड़े घटकों को छोटे घटकों में काट देते हैं: स्टार्च को
शर्कराओं में, माँस और \omterm{def:g2:animals-grow-up:born}{अंडों} के शरीर बनाने वाले पदार्थ को उसकी छोटी
इकाइयों में, चिकनाइयों को उनकी इकाइयों में। हर रस की अपनी ख़ासियतें हैं,
और हर रस अपने ही ठिकाने की \omterm{def:g6:exploring-our-environment:conditions}{परिस्थितियों} में सबसे अच्छा काम करता है।
\end{proposition}
\admitted

\begin{example}[मेज़ पर उतारा हुआ रूप]\label{ex:g7:digestion-and-nutrients:bench}
वह चिर-परिचित प्रयोग \omterm{def:g4:journey-of-food:digestion}{पाचन} को शीशे में चलाता है। शरीर की गरमाहट पर पके
स्टार्च की लेई वाली दो नलियाँ: एक में ताज़ी \omterm{def:g4:journey-of-food:tract}{लार}, दूसरी में उबाली हुई
(बिगाड़ी हुई) \omterm{def:g4:journey-of-food:tract}{लार}। एक घंटे बाद स्टार्च की जाँच दूसरी नली को जस-का-तस
पाती है --- अब भी स्टार्च --- जबकि पहली नली की लेई शर्कराओं में बदल
चुकी है, ठीक \cref{ex:g4:journey-of-food:sweet} वाली मिठास, और वह भी
बिना किसी मुँह के। तोड़-फोड़ अकेले रस ने की --- और उबाली हुई नली दिखाती
है कि उसका काम करने वाला पदार्थ गरमी से नष्ट हो जाता है, जैसे कोई नाज़ुक
मशीन।
\end{example}

\begin{remark}[तोड़-फोड़ करनी ही क्यों है?]\label{rem:g7:digestion-and-nutrients:why}
भोजन के बड़े घटक आँत की दीवार पार उतने ही नहीं कर सकते जितना कोई अलमारी
चिट्ठी वाले झिरी-दरवाज़े से नहीं गुज़र सकती। और शरीर उन्हें साबुत चाहता
भी नहीं: उसे \emph{पुर्ज़े} चाहिए --- वे छोटी इकाइयाँ, जिनसे वह अपना
ही पदार्थ दोबारा बना सके
(\cref{prop:g6:matter-from-food:transform})। \omterm{def:g4:journey-of-food:digestion}{पाचन} ढुलाई के लिए
\emph{और} दोबारा जोड़ने के लिए की गई तोड़-फोड़ है: हर खाने पर फ़र्नीचर
से खुले पुर्ज़ों का बंडल।
\end{remark}

\section{पोषक तत्व}

\begin{definition}[पोषक तत्वों के कुल]\label{def:g7:digestion-and-nutrients:families}
तोड़-फोड़ से \emph{पोषक तत्व}\index{पोषक तत्व} निकलते हैं --- यानी वे
छोटे, \omterm{prop:g4:journey-of-food:absorb}{रक्त} के लिए तैयार घटक:
\begin{itemize}
  \item \emph{शर्कराएँ} (सबसे बढ़कर \omterm{prop:g7:organs-at-work:bill}{ग्लूकोज़},
        \cref{prop:g7:organs-at-work:bill} वाला ईंधन), जो स्टार्च
        वाले और मीठे भोजनों से आती हैं;
  \item \emph{शरीर बनाने वाले पदार्थ} की छोटी इकाइयाँ, जो माँस,
        मछली, \omterm{def:g2:animals-grow-up:born}{अंडों}, \omterm{def:g2:animals-grow-up:born}{दूध} की चीज़ों और दालों से आती हैं --- यानी
        \cref{prop:g3:balanced-diet:jobs} वाली उसारी की ईंटें;
  \item \emph{चिकनाइयों} की इकाइयाँ --- गाढ़ा ईंधन भी और बनाने का
        सामान भी;
  \item और साथ में वे मुसाफ़िर जिन्हें कोई तोड़-फोड़ चाहिए ही नहीं:
        \emph{विटामिन}, \emph{\omterm{prop:g4:what-plants-need:needs}{खनिज लवण}}, और ख़ुद पानी।
\end{itemize}
\cref{def:g3:balanced-diet:families} वाले \omterm{def:g3:balanced-diet:families}{आहार-वर्ग} दरअसल इन्हीं पोषक
कुलों के आपूर्तिकर्ता हैं --- यानी सौदे की सूची के पीछे वाली सौदे की
सूची।
\end{definition}

\begin{example}[एक दोपहर का खाना, मद-दर-मद]\label{ex:g7:digestion-and-nutrients:lunch}
\cref{ex:g3:balanced-diet:lunch} वाला संतुलित खाना, पोषक तत्वों के रूप
में दोबारा पढ़ा हुआ: पास्ता --- स्टार्च, जो लीटरों \omterm{prop:g7:organs-at-work:bill}{ग्लूकोज़} में तोड़ा
जाता है; मछली --- शरीर बनाने वाला पदार्थ, अपनी इकाइयों तक कटा हुआ;
फलियों पर पड़ा तेल --- चिकनाई की इकाइयाँ; और फलियाँ तथा सेब --- विटामिन,
खनिज, पानी, और साथ में रेशा। ग़ौर करो, रेशा वह अपवाद है जो पूरे तंत्र को
साबित करता है: हमारा कोई रस उसे काट नहीं सकता, इसलिए वह कुछ भी पार नहीं
करता और आगे बढ़ जाता है --- बचे-खुचे का ढेर
(\cref{prop:g4:journey-of-food:absorb}), जो अपनी बुहारी ख़ुद लगाता चलता
है।
\end{example}

\section{वह बड़ी सरहद-पार}

\begin{proposition}[रसांकुरों पर अवशोषण]\label{prop:g7:digestion-and-nutrients:absorption}
\cref{prop:g4:journey-of-food:absorb} में वादा की गई वह सरहद-पार \omterm{def:g4:journey-of-food:tract}{छोटी
आँत} की भीतरी दीवार पर होती है, और दीवार उसी के लिए बनी है: तह पर तह, और
ऊपर से लाखों \emph{रसांकुरों}\index{रसांकुर} का क़ालीन --- यानी उँगली
जैसे उभार, एक मिलीमीटर ऊँचे, जिनमें से हर एक अपनी ही महीन रक्त-वाहिकाएँ
लपेटे रहता है। सरहद-पार वाली सतह अपने कमरे-भर-क़ालीन जितने आकार तक
\omterm{prop:g7:digestion-and-nutrients:absorption}{रसांकुरों} की वजह से ही पहुँचती है
(\cref{rem:g4:journey-of-food:surface}): तहों पर तहें, और उन पर
उँगलियाँ। \omterm{def:g7:digestion-and-nutrients:families}{पोषक तत्व} उनकी पतली दीवारों से पार होकर \omterm{prop:g4:journey-of-food:absorb}{रक्त} में उतरते हैं,
और \omterm{prop:g4:journey-of-food:absorb}{रक्त} उन्हें आम बँटवारे से पहले छाँटने-और-जमा करने वाले एक बड़े अंग
तक ले जाता है --- \emph{यकृत}\index{यकृत} तक।
\end{proposition}
\admitted

\begin{omfigure}
\begin{tikzpicture}[scale=0.9]
  % intestine wall with villi
  \draw[thick] (0,0) -- (0,3.4);
  \draw[thick] (0,0) -- (8.2,0);
  \node[font=\small, rotate=90] at (-0.4,1.7) {आँत की दीवार};
  % villi
  \foreach \x in {0.9,2.4,3.9,5.4,6.9} {
    \draw[thick, fill=omDef!8] (\x,0)
      .. controls ({\x-0.42},1.4) and ({\x-0.42},2.4) .. (\x,2.6)
      .. controls ({\x+0.42},2.4) and ({\x+0.42},1.4) .. (\x,0);
    % vessels inside
    \draw[omThm, thick] ({\x-0.14},0.15)
      .. controls ({\x-0.18},1.3) .. (\x,2.25)
      .. controls ({\x+0.18},1.3) .. ({\x+0.14},0.15);
  }
  \node[font=\footnotesize, align=center] at (4.1,-0.6)
    {हर रसांकुर अपनी ही महीन रक्त-वाहिकाएँ लपेटे है};
  % nutrients crossing
  \foreach \p in {(1.7,2.9),(3.2,3.0),(4.7,2.85),(6.2,2.95)} {
    \fill[omMeth] \p circle (0.07);
  }
  \node[font=\footnotesize, omMeth] at (6.9,3.25)
    {नली में पोषक तत्व};
  \draw[->, omMeth, thick] (3.2,2.9) -- (3.65,2.2);
  \node[font=\footnotesize, align=left] at (10.0,1.7)
    {सरहद-पार:\\ रसांकुर की\\ पतली दीवार से होकर,\\ रक्त में};
\end{tikzpicture}

{\small आँत का भीतरी क़ालीन: लाखों \omterm{prop:g7:digestion-and-nutrients:absorption}{रसांकुर}, और हर एक रक्त-वाहिकाओं से
भरी पतली दीवार वाली उँगली। सतह पर तह की हुई सतह --- यानी शरीर की वह
बड़ी सरहद-पार।}
\end{omfigure}

\begin{example}[वही विशेष-सूची, फिर से]\label{ex:g7:digestion-and-nutrients:spec}
बड़ी, पतली, नम, और \omterm{prop:g4:journey-of-food:absorb}{रक्त} से अस्तर की हुई: \omterm{prop:g7:digestion-and-nutrients:absorption}{रसांकुर} ठीक
\cref{prop:g7:how-animals-breathe:spec} वाली विशेष-सूची पर खरे उतरते
हैं --- बस \omterm{def:g7:what-is-respiration:respiration}{ऑक्सीजन} की जगह पोषक तत्वों के लिए। फेफड़ों की कोठरियाँ,
\omterm{def:g7:how-animals-breathe:four}{क्लोमों} के \omterm{def:g1:animals-around-us:covering}{पंख}, आँत के \omterm{prop:g7:digestion-and-nutrients:absorption}{रसांकुर}: शरीर को जहाँ भी किसी सरहद के पार बहुत
कुछ पहुँचाना हो, वह वही मशीन बनाता है। विशेष-सूची एक बार पहचान लो, फिर
वह हर जगह पहचानी जाती है।
\end{example}

\begin{method}[किसी पोषक तत्व का सिरे से सिरे तक पीछा करना]\label{met:g7:digestion-and-nutrients:trace}
किसी भी कौर के लिए एक \omterm{def:g7:digestion-and-nutrients:families}{पोषक तत्व} का पीछा करो:
\begin{enumerate}
  \item भोजन के वर्ग और उसके मुख्य घटकों के नाम बताओ;
  \item ठिकाने-दर-ठिकाने लिखो कि उसकी तोड़-फोड़ कहाँ शुरू और कहाँ
        ख़त्म होती है, और उस पर किन रसों ने काम किया;
  \item उसके साथ-साथ \omterm{prop:g7:digestion-and-nutrients:absorption}{रसांकुरों} पर सरहद पार करो, \omterm{prop:g7:digestion-and-nutrients:absorption}{यकृत} होते हुए \omterm{prop:g4:journey-of-food:absorb}{रक्त}
        में;
  \item और उसे पहुँचाओ: किसी काम करती \omterm{def:g5:first-look-at-cells:cell}{कोशिका} तक ईंधन के रूप में
        (\cref{prop:g7:organs-at-work:bill}), या किसी उसारी की जगह
        तक सामान के रूप में
        (\cref{def:g6:matter-from-food:production})।
\end{enumerate}
\end{method}

\begin{example}[रोटी पर चलाया हुआ पीछा]\label{ex:g7:digestion-and-nutrients:breadtrace}
रोटी: स्टार्च वाला वर्ग, और मुख्य रूप से स्टार्च। तोड़-फोड़ मुँह में
शुरू होती है (\omterm{def:g4:journey-of-food:tract}{लार} --- वही देर तक चबाने वाली मिठास), \omterm{def:g4:journey-of-food:tract}{आमाशय} के अम्ल में
ठहर जाती है, और \omterm{def:g4:journey-of-food:tract}{छोटी आँत} के रसों में पूरी होती है: \omterm{prop:g7:organs-at-work:bill}{ग्लूकोज़}। सरहद-पार:
\omterm{prop:g7:digestion-and-nutrients:absorption}{रसांकुर}, \omterm{prop:g4:journey-of-food:absorb}{रक्त}, \omterm{prop:g7:digestion-and-nutrients:absorption}{यकृत} --- जो उस \omterm{prop:g7:organs-at-work:bill}{ग्लूकोज़} का एक हिस्सा बाद के लिए जमा कर
लेता है, यानी अपनी बहुत-सी सेवाओं में से एक। पहुँचाना: दौड़ के बीच जाँघ
की एक \omterm{def:g3:bones-and-muscles:muscle}{पेशी} उसे उसी घंटे जला डालती है। चौथी कक्षा वाला कौर, पूरी तरह
हिसाब में।
\end{example}

\section{अभ्यास}

\begin{exercise}[$\star$]\label{exo:g7:digestion-and-nutrients:1}
\omterm{def:g4:journey-of-food:digestion}{पाचक रस} वह क्या करते हैं जो दाँत और मथना नहीं कर सकते?
\end{exercise}

\begin{exercise}[$\star$]\label{exo:g7:digestion-and-nutrients:2}
\omterm{def:g4:journey-of-food:tract}{लार} वाला दो-नलियों का प्रयोग बताओ: तैयारी, नतीजा, और दोनों निष्कर्ष।
\end{exercise}

\begin{exercise}[$\star$]\label{exo:g7:digestion-and-nutrients:3}
पोषक तत्वों के कुल गिनाओ, और हर एक का एक-एक भोजन-स्रोत भी।
\end{exercise}

\begin{exercise}[$\star$]\label{exo:g7:digestion-and-nutrients:4}
भोजन के किन घटकों को कोई तोड़-फोड़ नहीं चाहिए? और कौन-सा घटक हमारे रसों
को पूरी तरह हरा देता है --- और वह भी किस उपयोगी मक़सद से?
\end{exercise}

\begin{exercise}[$\star$]\label{exo:g7:digestion-and-nutrients:5}
\omterm{prop:g7:digestion-and-nutrients:absorption}{रसांकुर} क्या है? हर एक क्या लपेटे रहता है?
\end{exercise}

\begin{exercise}[$\star$]\label{exo:g7:digestion-and-nutrients:6}
आँत से आया \omterm{prop:g4:journey-of-food:absorb}{रक्त} आम बँटवारे से पहले कहाँ जाता है, और उस अंग की एक सेवा
का नाम भी बताओ।
\end{exercise}

\begin{exercise}[$\star\star$]\label{exo:g7:digestion-and-nutrients:7}
\cref{rem:g7:digestion-and-nutrients:why} वाली खुले पुर्ज़ों के बंडल
वाली तस्वीर समझाओ: भोजन को तोड़ना ही क्यों पड़ता है, दोनों कारण दो।
\end{exercise}

\begin{exercise}[$\star\star$]\label{exo:g7:digestion-and-nutrients:8}
उबाली हुई \omterm{def:g4:journey-of-food:tract}{लार} वाली नली क्यों मायने रखती है? उस तरीक़े का नाम बताओ, और
यह भी कि उबालना रस के काम करने वाले पदार्थ के बारे में क्या दिखाता है।
\end{exercise}

\begin{exercise}[$\star\star$]\label{exo:g7:digestion-and-nutrients:9}
\omterm{prop:g7:digestion-and-nutrients:absorption}{रसांकुरों} को फेफड़ों की कोठरियों और \omterm{def:g7:how-animals-breathe:four}{क्लोमों} के \omterm{def:g1:animals-around-us:covering}{पंखों} के बग़ल में रखो:
साझी विशेष-सूची बताओ, और हर मशीन का माल भी।
\end{exercise}

\begin{exercise}[$\star\star$]\label{exo:g7:digestion-and-nutrients:10}
\cref{met:g7:digestion-and-nutrients:trace} को एक उबले \omterm{def:g2:animals-grow-up:born}{अंडे} (शरीर
बनाने वाला पदार्थ) पर चलाओ, और वह भी भरती हुई खरोंच की मरम्मत की जगह
तक (\cref{exo:g6:cell-unit-of-life:10})।
\end{exercise}

\begin{exercise}[$\star\star$]\label{exo:g7:digestion-and-nutrients:11}
एक आनन-फ़ानन वाला परहेज़ हफ़्तों तक शरीर बनाने वाले सारे भोजन छोड़ देता
है। \cref{def:g7:digestion-and-nutrients:families} और
\cref{def:g6:matter-from-food:production} की मदद से भविष्यवाणी करो कि
अब भी बढ़ते हुए शरीर में क्या कम पड़ेगा और सबसे पहले क्या भुगतेगा।
\end{exercise}

\begin{exercise}[$\star\star\star$]\label{exo:g7:digestion-and-nutrients:12}
``\omterm{def:g4:journey-of-food:digestion}{पाचन} किसी खाने को पहुँचाने की सूची में बदल देता है।'' इस वाक्य का
सिरे से सिरे तक बचाव करो --- रस, \omterm{def:g7:digestion-and-nutrients:families}{पोषक तत्व}, \omterm{prop:g7:digestion-and-nutrients:absorption}{रसांकुर}, \omterm{prop:g7:digestion-and-nutrients:absorption}{यकृत}, पहुँचाना
--- और वह भी एक ही अनुच्छेद में।
\end{exercise}

\section{समस्या: शीशे की आँत}

\begin{problem}[{सप्ताहांत समस्या --- मेज़ पर, नली-दर-नली पाचन
बनाना}]\label{pb:g7:digestion-and-nutrients:1}
कक्षा एक ``शीशे की आँत'' बनाती है: शरीर की गरमाहट पर रखी नलियों की एक
चौकी, जिसमें हर नली नली-तंत्र के एक ठिकाने की नक़ल करती है। नली म
(मुँह): स्टार्च की लेई और ताज़ी \omterm{def:g4:journey-of-food:tract}{लार}। नली म$'$: स्टार्च की लेई और उबाली
हुई \omterm{def:g4:journey-of-food:tract}{लार}। नली प (पेट): \omterm{def:g2:animals-grow-up:born}{अंडे} की सफ़ेदी के टुकड़े और दवा वाले से लिया
\omterm{def:g4:journey-of-food:tract}{आमाशय} का रस। नली प$'$: वही टुकड़े, सादे पानी में। सब दो घंटे चलती हैं;
फिर जाँचें --- स्टार्च की, शर्कराओं की, और टुकड़ों पर एक नज़र।

\medskip\noindent\textbf{भाग I --- भविष्यवाणियाँ और नतीजे।}
\begin{enumerate}
  \item चारों नलियों के दो घंटे बाद के नतीजों की भविष्यवाणी करो।
  \item नली म मीठी हो जाती है, म$'$ नहीं। निष्कर्ष बताओ, और म$'$ की
        ठीक-ठीक भूमिका भी।
  \item नली प के टुकड़े सिकुड़ते, मुलायम होते और साफ़ हो जाते हैं;
        प$'$ के जस-के-तस बैठे रहते हैं। \omterm{def:g4:journey-of-food:tract}{आमाशय} के रस ने क्या किया ---
        और पोषक तत्वों का कौन-सा कुल बन रहा है?
  \item किसी भी नली में न चबाना हुआ, न मथना। यह चौकी इस बारे में क्या
        साबित करती है कि \omterm{def:g4:journey-of-food:digestion}{पाचन} का असली काम कहाँ है?
\end{enumerate}

\medskip\noindent\textbf{भाग II --- नमूने को बेहतर बनाना।}
\begin{enumerate}[resume]
  \item एक विद्यार्थी ``सफ़ाई के लिहाज़ से'' नली म को बर्फ़ीले
        डिब्बे के तापमान पर चलाने का प्रस्ताव रखता है।
        \cref{prop:g6:decomposers-and-soil:speed} वाला ढर्रा रसों पर
        लगाकर एतराज़ करो।
  \item एक और विद्यार्थी दोनों रस मिलाकर एक ही बड़ी नली का प्रस्ताव
        रखता है। यह मिलावट
        \cref{prop:g7:digestion-and-nutrients:chemical} के किस
        ठिकाना-दर-ठिकाना वाले तथ्य को धुँधला कर देगी?
  \item शीशे की आँत की कोई दीवार नहीं: कुछ भी अवशोषित नहीं होता। पूरे
        नमूने को आगे कौन-सी रचना चाहिए होगी, और उसके सामान में कौन-से
        दो गुण होने ही चाहिए?
  \item पाँचवीं नली में हर मौजूद रस के साथ रेशा --- चोकर --- डालो।
        दो घंटे बाद उसकी हालत की भविष्यवाणी करो, और नमूने के लिए
        उसका मतलब भी।
\end{enumerate}

\medskip\noindent\textbf{भाग III --- शीशे से आँत तक।}
\begin{enumerate}[resume]
  \item शरीर में नली म की शर्कराएँ आगे \omterm{prop:g7:digestion-and-nutrients:absorption}{रसांकुरों} से मिलतीं। उनकी
        यात्रा चार पड़ावों में आगे बढ़ाओ।
  \item \omterm{def:g4:journey-of-food:tract}{आमाशय} का रस कड़े अम्ल में काम करता है; और आँत के रसों को वह
        अम्ल बेअसर किया हुआ चाहिए। यह ``हर रस अपने ही ठिकाने की
        \omterm{def:g6:exploring-our-environment:conditions}{परिस्थितियों} में'' वाली बात के बारे में क्या कहता है --- और
        ठिकानों को क्रम से निगलने के बारे में क्या?
  \item एक दवा कंपनी उन लोगों के लिए रस के कैप्सूल बेचती है जिनके
        अपने ठिकाने कम काम करते हैं। चौकी की मदद से समझाओ कि ऐसा
        इलाज मुमकिन ही क्यों है।
  \item इस प्रयोग का सार एक वाक्य में दो, और उसमें \emph{रस},
        \emph{\omterm{def:g7:digestion-and-nutrients:families}{पोषक तत्व}} तथा \emph{नियंत्रण-नलियाँ} शब्द रखो।
\end{enumerate}
\end{problem}
